\documentclass[12pt,a4paper]{article}
\usepackage[margin=2.5cm]{geometry}
\usepackage[T1]{fontenc}
\usepackage[utf8]{inputenc}
\usepackage[brazil]{babel}
\usepackage{lmodern}
\usepackage{hyperref}
\usepackage{enumitem}

\title{Sistema Instituto Alento \\ Guia para Cliente}
\author{Instituto Alento}
\date{\today}

\begin{document}
\maketitle
\tableofcontents
\newpage

\section{Objetivo do sistema}
O Sistema Instituto Alento organiza o atendimento social e de saude em um unico lugar.
Ele ajuda a equipe a cadastrar familias, acompanhar dependentes, marcar consultas, controlar presenca e registrar historico.

\section{Perfis de acesso}
\begin{itemize}[leftmargin=*]
\item \textbf{Admin}: configura o sistema e acompanha toda a operacao.
\item \textbf{Medico/Voluntario}: organiza agenda e registra atendimentos.
\item \textbf{Familia}: acompanha dependentes, consultas e notificacoes.
\end{itemize}

\section{Principais recursos}
\begin{itemize}[leftmargin=*]
\item Dashboard com indicadores da operacao.
\item Cadastro de familias e dependentes.
\item Agenda de consultas e atendimentos.
\item Registro de presenca, falta e justificativa.
\item Modulos integrados (Help Desk / HDI), quando liberados.
\end{itemize}

\section{Fluxo simples do dia a dia}
\subsection{1) Cadastrar familia}
Registrar dados basicos do responsavel e dos dependentes.

\subsection{2) Organizar atendimento}
Selecionar profissional, data, hora e sala (quando necessario).

\subsection{3) Realizar atendimento}
No dia do compromisso, confirmar presenca e registrar observacoes.

\subsection{4) Acompanhar resultado}
Usar painel e agenda para ver pendencias, historico e proximos passos.

\section{Beneficios para a instituicao}
\begin{itemize}[leftmargin=*]
\item Menos retrabalho e menos perda de informacao.
\item Mais rapidez para localizar dados de familias e dependentes.
\item Visao clara do que foi feito e do que esta pendente.
\item Base mais forte para decisoes da equipe gestora.
\end{itemize}

\section{Boas praticas para a equipe}
\begin{itemize}[leftmargin=*]
\item Manter dados sempre atualizados.
\item Registrar presenca no mesmo dia do atendimento.
\item Usar filtros rapidos para priorizar casos.
\item Evitar compartilhar senha de usuario.
\end{itemize}

\section{Canais de suporte}
Em caso de duvida:
\begin{itemize}[leftmargin=*]
\item acionar equipe interna de administracao;
\item abrir chamado no modulo de suporte, quando disponivel.
\end{itemize}

\end{document}
